\cotefile{27}
\scanpage{15}

$1^\circ)$ On peut prendre pour $n$ le degré de \unsure{tte|toute} extension
finie $k'$ de $k$ ayant les propriétés suivantes

a) $G_{k'\,\mathrm{réd}}$ est lisse sur $k'$ i.e. séparable $/k'$

b) \deleted{$G_{k'\,\mathrm{réd}}$ admet un Borel}
les pts de $G/G^0$ sont \deleted{définis}\inserted{rationnels} \uline{sur $k'$}.

c) \inserted{``le'' radical $R'$} \unsure{de} $G^0_{k'\,\mathrm{réd}}$ \deleted{est} \inserted{est défini sur $k'$, et splitté}
\deleted{admet un Borel $B'$ qui est splitté} sur $k'$ [suite de composition à
quotients successifs $G_a$ ou $G_m$]

d) \inserted{le \unsure{gpe rédt|groupe rédt|hyp} $G^0_{k'\,\mathrm{réd}}/R'$ est \deleted{un groupe \unsure{réductif|$k$ simple}} déployé}
\deleted{\unsure{de|le} $T'$ est un tore maximal de $B'$}
\deleted{[nécessairement splitté] et $N(T')$, $C(T')$}
\deleted{\unsure{leur|le|son} normalisateur resp\textsuperscript{t} centralisateur, dans}

\deleted{Par suite,}

\deleted{$1^\circ)$ Si $G$ est unipotent \inserted{connexe}, on peut prendre
$n$ une puissance de la caractéristique}

\deleted{$2^\circ)$ Si $G$ est \deleted{réductif} $\frac12$ simple, alors
($\operatorname{Hom}(G,G_m)$ étant nul, après \deleted{ext} \unsure{ext.} du
corps de base] $\operatorname{Pic}$}
\deleted{$G$ est lisse}

$2^\circ)$ Supposons que $\operatorname{Hom}_{\unsure{\bar k\text{-}\mathrm{gr}|\bar k\text{-}\mathrm{gp}}}(G_{\bar k},G_{m\,\bar k})=0$,
et que $G$ soit \inserted{lisse et} connexe. Soit $\widetilde k$ la
clôture \inserted{algébrique} \uuline{séparable} de $k$. Alors $\operatorname{Pic}(G)\simeq\operatorname{Pic}(G_{\widetilde k})$

\deleted{En particulier, si $\operatorname{Hom}(G_{\unsure{k|\bar k}},G_m)$}

Si \inserted{on a} $\underline{\operatorname{Hom}}_{\unsure{\widetilde k\text{-}\mathrm{gr}|\widetilde k\text{-}\mathrm{gp}}}(G_{\widetilde k},G_{m\,k})=0$ [ce qui est
vrai si en plus de la \unsure{relation|condition} \unsure{précédte|précédente} sur $\bar k$,
$G$ est de \uline{rang unipotent} \uline{nul}] on a $=\operatorname{Pic}(G)\simeq\operatorname{Pic}(G_{\widetilde k})$

% archive sheet number: 14
